\stfhIndex{I199}{Yebaomianjitongpei}{野保面积同胚（wild area-preserving homeomorphism）}
\section{\texorpdfstring{$D^2$}{D2} 上的野保面积同胚}

本节描述 $\operatorname{Sympeo}(D^2,\partial D^2)$ 中一个紧支保面积同胚的例子；若
猜想~\ref{conj:22-calabi-homeomorphism-extension} 成立，它将不属于
$\Hameo(D^2,\partial D^2)$。由定理 6.4.1，这将蕴含
$\Hameo(D^2,\partial D^2)$ 是
$\operatorname{Sympeo}(D^2,\partial D^2)$ 的真正规子群。再结合 (Oh06b), (Sik07)
中的下述光滑化定理，上述一系列陈述将导出
$\operatorname{Homeo}_\Omega(D^2,\partial D^2)$ 非单纯。

\begin{theorem}
\label{thm:22-area-preserving-smoothing}
对 $D^2$ 上的标准面积形式 $\Omega$，把它也看作辛形式 $\omega=\Omega$，有
\[
 \operatorname{Sympeo}(D^2,\partial D^2)
 =\operatorname{Homeo}_\Omega(D^2,\partial D^2).
\]
\end{theorem}

下面集中构造 $D^2$ 上野保面积同胚的例子。为此，需要考虑 $D^2$ 的重标度映射
\[
 R_a:D^2(1)\longrightarrow D^2(a)\subset D^2(1),
 \qquad 0<a<1,
\]
在 $\Hameo(D^2,\partial D^2)$ 上的共轭作用；这里 $D^2(a)$ 是以原点为中心、半径为
$a$ 的圆盘。$R_a$ 是共形辛映射，所以只要其共轭作用有定义，就把辛映射映成辛映射。

此外，右复合 $R_a$ 定义映射
\[
 \phi\longmapsto\phi\circ R_a:\Hameo(D^2,\partial D^2)
 \longrightarrow\Hameo(D^2(a),\partial D^2(a))
 \subset\Hameo(D^2,\partial D^2),
\]
而复合 $R_a^{-1}$ 定义映射
\[
 \Hameo(D^2(a),\partial D^2(a))
 \longrightarrow\Hameo(D^2,\partial D^2).
\]
关于亚历山大同痕下卡拉比不变量的行为，有如下重要公式。

\begin{lemma}
\label{lem:22-calabi-rescaling}
令 $\lambda$ 为 $D^2$ 上给定的拓扑哈密顿道路，并定义映射 $\lambda_a$ 为
\[
 \lambda_a(t,x)=
 \begin{cases}
  a\lambda(t,x/a),&|x|\leq a(1-\eta),\\
  x,&\text{其他情形},
 \end{cases}
 \qquad 0<a\leq1.
\]
则 $\lambda_a$ 也是 $D^2$ 上的拓扑哈密顿道路，并满足
\begin{equation}
 \overline{\operatorname{Cal}}^{\mathrm{path}}(\lambda_a)
 =a^4\overline{\operatorname{Cal}}^{\mathrm{path}}(\lambda).
 \label{eq:22-calabi-rescaling}
\end{equation}
\end{lemma}

\begin{proof}
直接计算可知，$\lambda_a$ 由如下（唯一）拓扑哈密顿量生成：
\[
 \operatorname{Dev}(\lambda_a)(t,x)=
 \begin{cases}
  a^2H(t,x/a),&|x|\leq a(1-\eta),\\
  0,&\text{其他情形},
 \end{cases}
\]
其中 $H=\operatorname{Dev}(\lambda)$。显然，若 $\lambda_i$ 是 $\lambda$ 的一列光滑
哈密顿逼近，且 $\lambda_{i,a}$ 由与 $\lambda_i$ 相同的公式定义，则右端函数是
$\operatorname{Dev}(\lambda_{i,a})$ 的哈密顿极限。由此推出
\begin{align*}
 \overline{\operatorname{Cal}}^{\mathrm{path}}(\lambda_a)
 &=\int_0^1\int_{D^2(a(1-\eta))}a^2H\!\left(t,\frac xa\right)
   \Omega\wedge dt\\
 &=a^4\int_0^1\int_{D^2}H(t,y)\Omega\wedge dt
 =a^4\overline{\operatorname{Cal}}^{\mathrm{path}}(\lambda).
\end{align*}
这证明了~\eqref{eq:22-calabi-rescaling}。
\end{proof}

下面构造一个野保面积同胚的例子 (Oh10)。

\begin{example}
\label{ex:22-wild-area-preserving-homeomorphism}
作好上述准备，考虑二进数集合 $1/2^k$（$k=0,\ldots$）。令 $(r,\theta)$ 为 $D^2$
上的极坐标，则标准面积形式为
\[
 \omega=r\,dr\wedge d\theta.
\]
考虑如下形式的映射 $\phi_k:D^2\to D^2$：
\[
 \phi_k=\phi_{\rho_k}:(r,\theta)\longmapsto(r,\theta+\rho_k(r)),
\]
其中 $\rho_k:(0,1]\to[0,\infty)$ 是支集在 $(0,1)$ 中的光滑函数。于是
$\phi_{\rho_k}$ 是由自治哈密顿量
\[
 F_{\phi_k}(r,\theta)=-\int_1^r s\rho_k(s)\,ds
\]
生成的保面积映射。因此其卡拉比不变量为
\begin{equation}
 \operatorname{Cal}(\phi_k)
 =-\int_{D^2}\left(\int_1^r s\rho_k(s)\,ds\right)r\,dr\,d\theta
 =\pi\int_0^1r^3\rho_k(r)\,dr.
 \label{eq:22-twist-calabi}
\end{equation}
\footnote{译者注：原书式 (22.7.75) 的末积分变量印作 $dt$，但被积函数及积分上下限
均以 $r$ 为变量；译文改作 $dr$。}
如下选取 $\rho_k$：
\begin{enumerate}[label=(\arabic*)]
\item $\rho_k$ 的支集在 $1/2^k<r<1/2^{k-1}$ 中；
\item 对每个 $k=1,\ldots$，当 $r\in(1/2^k,1/2^{k-1})$ 时，
\begin{equation}
 \rho_k(r)=2^4\rho_{k-1}(2r);
 \label{eq:22-twist-rescaling}
\end{equation}
\item $\operatorname{Cal}(\phi_1)=1$。
\end{enumerate}

按构造，$\phi_k$ 的支集彼此不交，故可以自由复合而无需考虑复合次序。因此无穷乘积
\[
 \prod_{k=0}^\infty\phi_k
\]
\footnote{译者注：原书这里把乘积下标写成 $k=0$，但紧接着的分段定义、后文
$\psi_N$ 与定理证明均从 $k=1$ 开始；并且按性质 (1)，$\rho_0$ 在其定义域
$(0,1]$ 内不可能支撑于 $1<r<2$。这里疑应为 $\prod_{k=1}^\infty\phi_k$；译文保留
原下标并提示索引不一致。}
良定，并定义一个除原点外光滑的连续映射；在原点处，$\phi_\rho$ 连续但不可微。
这个无穷乘积也可写成同胚 $\phi_\rho$，其取值为 $\phi_\rho(0)=0$ 以及
\[
 \phi_\rho(r,\theta)=(r,\theta+\rho(r)),
\]
其中光滑函数 $\rho:(0,1]\to\mathbb R$ 定义为
\[
 \rho(r)=\rho_k(r),\qquad
 r\in[1/2^k,1/2^{k-1}],\quad k=1,2,\ldots.
\]
容易验证 $\phi_\rho$ 在 $D^2\setminus\{0\}$ 上光滑，在 $0$ 处也连续，并与上述无穷
乘积一致。显然 $\phi_{-\rho}$ 是 $\phi_\rho$ 的逆，故它是同胚。此外，
\[
 \phi_\rho^*(r\,dr\wedge d\theta)=r\,dr\wedge d\theta
 \qquad\text{于 }D^2\setminus\{0\},
\]
所以 $\phi_\rho$ 确实保面积。
\end{example}

下述引理将在定理~\ref{thm:22-wild-not-hameomorphism} 的证明中起重要作用。

\begin{lemma}
\label{lem:22-twist-self-similarity}
令 $\phi_k$ 为例~\ref{ex:22-wild-area-preserving-homeomorphism} 中的微分同胚。则
\begin{equation}
 R_{\frac12}\circ\phi_{k-1}^{2^4}\circ R_{\frac12}^{-1}=\phi_k.
 \label{eq:22-twist-conjugacy}
\end{equation}
特别地，
\begin{equation}
 \operatorname{Cal}(\phi_k)=\operatorname{Cal}(\phi_{k-1}).
 \label{eq:22-twist-calabi-constant}
\end{equation}
\end{lemma}

\begin{proof}
由~\eqref{eq:22-twist-rescaling}，计算得
\[
 R_{\frac12}\circ\phi_{k-1}\circ R_{\frac12}^{-1}(r,\theta)
 =(r,\theta+\rho_{k-1}(2r))
 =\left(r,\theta+\frac1{2^4}\rho_k(r)\right).
\]
迭代该恒等式 $2^4$ 次，便由~\eqref{eq:22-twist-rescaling} 得到
\eqref{eq:22-twist-conjugacy}。再由它与~\eqref{eq:22-calabi-rescaling} 得到
\eqref{eq:22-twist-calabi-constant}。
\end{proof}

该引理与式~\eqref{eq:22-twist-rescaling} 立即给出如下推论。
\footnote{译者注：推论实际还使用例~\ref{ex:22-wild-area-preserving-homeomorphism}
中的初值 $\operatorname{Cal}(\phi_1)=1$，且递推恒等式直接来自
式~\eqref{eq:22-twist-calabi-constant}；原书这里只引用式 (22.7.76)。}

\begin{corollary}
\label{cor:22-twist-calabi-one}
对所有 $k=1,\ldots$，
\[
 \operatorname{Cal}(\phi_k)=1.
\]
\end{corollary}

现在可以证明下述定理。

\begin{theorem}
\label{thm:22-wild-not-hameomorphism}
若猜想~\ref{conj:22-calabi-homeomorphism-extension} 成立，则
$\phi_\rho$ 不属于 $\Hameo(D^2,\partial D^2)$。
\end{theorem}

\begin{proof}
反设 $\phi_\rho\in\Hameo(D^2,\partial D^2)$，即存在道路
$\lambda\in\mathcal P^{\mathrm{ham}}
(\operatorname{Sympeo}(D^2,\partial D^2),id)$，满足 $\lambda(1)=\phi_\rho$。
于是其卡拉比不变量有有限值，记作
\begin{equation}
 \overline{\operatorname{Cal}}(\phi_\rho)
 =\overline{\operatorname{Cal}}^{\mathrm{path}}(\lambda)=C_1
 \label{eq:22-wild-calabi-value}
\end{equation}
其中 $C_1\in\mathbb R$。

写 $\phi_\rho=\psi_N\widetilde\psi_N$，其中
\[
 \psi_N=\prod_{i=1}^N\phi_i,
 \qquad
 \widetilde\psi_N=\prod_{i=N+1}^\infty\phi_i.
\]
取 $N=1$，由 $\overline{\operatorname{Cal}}$ 的同态性质得到
\begin{equation}
 C_1=\overline{\operatorname{Cal}}(\psi_1)
 +\overline{\operatorname{Cal}}(\widetilde\psi_1).
 \label{eq:22-wild-calabi-split}
\end{equation}
注意 $\psi_N$ 光滑，故显然属于 $\Hameo(D^2,\partial D^2)$。若 $\phi_\rho$ 属于该群，
由 $\Hameo(D^2,\partial D^2)$ 的群性质，$\widetilde\psi_N$ 也属于该群。又由
推论~\ref{cor:22-twist-calabi-one}，
\[
 \overline{\operatorname{Cal}}(\psi_1)
 =\operatorname{Cal}(\phi_1)=1,
\]
所以
\begin{equation}
 \overline{\operatorname{Cal}}(\widetilde\psi_1)=C_1-1.
 \label{eq:22-wild-tail-calabi}
\end{equation}

另一方面，把~\eqref{eq:22-twist-conjugacy} 迭代应用于无穷乘积
\[
 \widetilde\psi_1=\prod_{i=2}^\infty\phi_i,
\]
可知 $\widetilde\psi_1$ 满足
\begin{equation}
 \widetilde\psi_1(r,\theta)=
 \begin{cases}
  R_{\frac12}\circ\phi_\rho^{2^4}\circ R_{\frac12}^{-1}(r,\theta),
       &0<r\leq\frac12,\\
  (r,\theta),&\frac12\leq r\leq1.
 \end{cases}
 \label{eq:22-wild-tail-rescaling}
\end{equation}
特别地，
\begin{equation}
 \widetilde\psi_1
 =\overline{\operatorname{ev}}_1((\lambda^{2^4})_{1/2}).
 \label{eq:22-wild-tail-evaluation}
\end{equation}
将引理~\ref{lem:22-calabi-rescaling} 用于 $a=1/2$，再用
$\overline{\operatorname{Cal}}^{\mathrm{path}}$ 的同态性质，得到
\begin{align}
 \overline{\operatorname{Cal}}(\widetilde\psi_1)
 &=\overline{\operatorname{Cal}}^{\mathrm{path}}((\lambda^{2^4})_{1/2})
 =\left(\frac12\right)^4
 \overline{\operatorname{Cal}}^{\mathrm{path}}(\lambda^{2^4})\notag\\
 &=\left(\frac12\right)^4 2^4
 \overline{\operatorname{Cal}}^{\mathrm{path}}(\lambda)
 =\overline{\operatorname{Cal}}^{\mathrm{path}}(\lambda)=C_1.
 \label{eq:22-wild-tail-calabi-rescaled}
\end{align}
显然，式~\eqref{eq:22-wild-tail-calabi} 与
\eqref{eq:22-wild-tail-calabi-rescaled} 相矛盾。
\end{proof}
